\subsection{Stokes' formula for pieces}

In this section, X denotes a manifold of class \(C^{r}\) ( \(r \geqslant 2\) ) pure of finite dimension n and separated. A denotes a piece of X and i the canonical injection of \(\partial A\) into X. E denotes a Banach space.

11.2.1. Let \(x \in \partial A\) and let \(\xi\) be an orientation of \(T_x(\partial A)\); one denotes \(\tilde{i}_x(\xi)\) the orientation of \(T_x(X)\) containing the elements \(v \wedge u\), where \(v\) is a strictly outward vector for A at x (11.1.4) and where \(u\) is a nonzero element of \(\bigwedge^{n-1} T_x(\partial A)\) belonging to the orientation \(\xi\). The map \(\tilde{i}_x\) is a bijection of \(\text{Or}(T_x(\partial A))\) onto \(\text{Or}(T_x(X))\). The maps \(\tilde{i}_x\), for \(x \in \partial A\), define a morphism \(\tilde{i} \colon \partial \widetilde{A} \to \widetilde{X}\) which is an orientation of \(i\) (10.2.5). If \(\xi\) is an orientation of X, the orientation of \(\partial A\) associated with \(\xi\) by \(\tilde{i}\) (10.2.5) is said to be defined by \(\xi\).

Example. --- When \(X = R^{n}\), \(\xi\) is the usual orientation of \(R^{n}\), and A is a closed ball of radius \(\rho>0\), the orientation of the sphere \(\partial A\) defined by \(\xi\) is the canonical orientation (10.2.8, b)).

11.2.2. Let \(\omega\) be a twisted differential form of degree \(p\) on X with values in a Banach space E (10.4.1); the inverse image \(i^{*}(\omega)\) of \(\omega\) by the oriented morphism \(i\colon\partial A\to X\) is denoted \(\omega|\partial A\) and is called the form induced by \(\omega\) on \(\partial A\) (cf. 10.4.3).

11.2.3. Let us make one of the following two assumptions:

\begin{enumerate}
\def\labelenumi{(\roman{enumi})}
\tightlist
\item
  \(\omega\) is a twisted differential form of degree \(n - 1\) on X, with values in E;
\item
  X is oriented, \(\partial A\) is equipped with the corresponding orientation (11.2.1), and \(\omega\) is a differential form of degree \(n - 1\) on X with values in E.
\end{enumerate}

Suppose moreover that \(\omega\) is of class \(C^{1}\) and that the intersection of A and the support of \(\omega\) is compact. The exterior differential \(d\omega\) of \(\omega\) is continuous (8.3.5 and 10.3.4).

The characteristic function of A is essentially integrable with respect to the vector measure defined by \(d\omega\) on X, and the differential form \(\omega|\partial\mathrm{A}\) of degree \(n-1\) on \(\partial\mathrm{A}\) is integrable (10.4.3 and 10.4.4). We have

\[
\int_ {\mathrm{A}} d \omega = \int_ {\partial \mathrm{A}} \omega \quad (\text{Stokes' formula})
\]

11.2.4. Let \(\alpha\) be a twisted differential form of degree \(n\) on X with values in a Banach space E, with compact support and of class \(\mathbf{C}^1\). For there to exist a twisted differential form \(\omega\) of degree \(n - 1\) on X, with values in E, with compact support and of class \(\mathbf{C}^1\), such that \(\alpha = d\omega\), it is necessary that \(\int_{\mathbb{X}} \alpha = 0\). If X is connected, this condition is sufficient; if moreover \(\alpha\) is of class \(\mathbf{C}^k\) (\(k\leqslant r-1,\ k\neq\omega\)), one can choose \(\omega\) of class \(\mathbf{C}^k\).

11.2.5. Suppose that X is the real manifold underlying \(\mathbf{C}\), that E is a complex Banach space, and that A is a compact piece of X. Endow X with the orientation defined by its complex structure (10.2.7). Let f be a continuous map from A into E whose restriction to the interior \(\mathring{A}\) of A is holomorphic. Denote by \(dz\) the differential of the injection \(z\colon\partial A\to\mathbf{C}\). The form \(f\cdot dz\), the product of f and dz, is a differential form of degree 1 on \(\partial A\), with values in E, of class \(C^{0}\), and one has:

\[
\int_ {\partial \mathbf {A}} f\, d z = 0 \quad (\text{Cauchy's formula})
\]

When \(f\) extends to a holomorphic map, still denoted \(f\), from an open neighborhood U of A into E, the differential form \(f.dz\) (where z now denotes the canonical injection of U into \(\mathbf{C}\)) is of class \(C^{\omega}\) on U and its exterior differential is zero.

11.2.6 (“Derivative of an integral”). Suppose X is oriented. Let Y be a manifold of class \(C^{r}\), and let \(\alpha\) be a differential form of degree \(n\) on Y, with values in E, of class \(C^{1}\), and with compact support. Let I be an open subset of R containing 0, and let \(g: I \times X \to Y\) be a morphism of class \(C^{r}\); for \(t \in I\), denote by \(g_{t}\) the map \(x \mapsto g(t, x)\) from X into Y. Denote by \(\psi\) the morphism from X into T(Y) defined by \(\psi(x) = T_{(0,x)}(g)(1,0)\); it is a lifting of \(g_{0}\) (8.6.5). Suppose that the restriction of \(\operatorname{pr}_{1}\colon I\times A\to I\) to the intersection of \(I\times A\) and the support of \(g^{*}(\alpha)\) is proper (TG, I, § 10). For every \(t \in I\), the intersection of A and the support of \(g_{t}^{*}(\alpha)\) is then compact; the map \(t \mapsto \int_{A} g_{t}^{*}(\alpha)\) is of class \(C^{1}\) on I and its derivative at the origin is given by the formula

\[
\frac {d}{d t} \left(\int_ {\mathrm{A}} g _ {t} ^ {*} (\alpha)\right) _ {t = 0} = \int_ {\mathrm{A}} \theta_ {\psi}. \alpha = \int_ {\mathrm{A}} i (\psi) d \alpha + \int_ {\partial \mathrm{A}} i (\psi) \alpha ,
\]

cf. no. 8.6.

